
\documentclass{article}
\usepackage{colortbl}
\usepackage{makecell}
\usepackage{multirow}
\usepackage{supertabular}

\begin{document}

\newcounter{utterance}

\centering \large Interaction Transcript for game `hot\_air\_balloon', experiment `air\_balloon\_survival\_en\_reasoning off\_hard', episode 4 with gpt{-}5{-}mini{-}2025{-}08{-}07{-}t1.0.
\vspace{24pt}

{ \footnotesize  \setcounter{utterance}{1}
\setlength{\tabcolsep}{0pt}
\begin{supertabular}{c@{$\;$}|p{.15\linewidth}@{}p{.15\linewidth}p{.15\linewidth}p{.15\linewidth}p{.15\linewidth}p{.15\linewidth}}
    \# & $\;$A & \multicolumn{4}{c}{Game Master} & $\;\:$B\\
    \hline

    \theutterance \stepcounter{utterance}  
    & & \multicolumn{4}{p{0.6\linewidth}}{
        \cellcolor[rgb]{0.9,0.9,0.9}{
            \makecell[{{p{\linewidth}}}]{
                \texttt{\tiny{[P1$\langle$GM]}}
                \texttt{Stai partecipando a un gioco di negoziazione collaborativa.} \\
\\ 
\texttt{Insieme a un altro partecipante, dovete accordarvi su un unico set di oggetti da conservare. Ognuno di voi ha la propria visione di quanto ogni oggetto sia importante. Non sai come l'altro partecipante valuta gli oggetti. Inoltre, ti viene fornito lo sforzo che ogni oggetto richiede.} \\
\texttt{Potete accordarvi su un set solo se lo sforzo totale degli oggetti selezionati non supera un limite condiviso:} \\
\\ 
\texttt{LIMITE = 2984} \\
\\ 
\texttt{Ecco i valori di sforzo dei singoli oggetti:} \\
\\ 
\texttt{Sforzo oggetto = \{"A21": 350, "B93": 276, "A11": 616, "B33": 737, "C08": 531, "B39": 389, "B96": 24, "C85": 125, "A82": 338, "B31": 356, "B20": 143, "A60": 117, "B09": 257, "C81": 921, "B21": 789\}} \\
\\ 
\texttt{Ecco la tua visione personale sull'importanza di ogni oggetto:} \\
\\ 
\texttt{Valori di importanza dell'oggetto = \{"A21": 138, "B93": 583, "A11": 868, "B33": 822, "C08": 783, "B39": 65, "B96": 262, "C85": 121, "A82": 508, "B31": 780, "B20": 461, "A60": 484, "B09": 668, "C81": 389, "B21": 808\}} \\
\\ 
\texttt{Obiettivo:} \\
\\ 
\texttt{Il tuo obiettivo è negoziare un set condiviso di oggetti che ti avvantaggi il più possibile (cioè, massimizzi l'importanza totale per TE), rimanendo entro il LIMITEE. Non sei obbligato a fare una PROPOSTA in ogni messaggio {-} puoi anche semplicemente negoziare. Tutte le tattiche sono permesse!} \\
\\ 
\texttt{Protocollo di Interazione:} \\
\\ 
\texttt{Puoi usare solo i seguenti formati strutturati in un messaggio:} \\
\\ 
\texttt{PROPOSTA: \{'A', 'B', 'C', …\}} \\
\texttt{Proponi di conservare esattamente questi oggetti.} \\
\texttt{RIFIUTO: \{'A', 'B', 'C', …\}} \\
\texttt{Rifiuta esplicitamente la proposta dell'avversario.} \\
\texttt{ARGOMENTAZIONE: \{'...'\}} \\
\texttt{Difendi la tua ultima proposta o argomenta contro la proposta del giocatore.} \\
\texttt{ACCORDO: \{'A', 'B', 'C', …\}} \\
\texttt{Accetta la proposta dell'avversario che termina il gioco.} \\
\\ 
\\ 
\texttt{Regole:} \\
\\ 
\texttt{Puoi ACCETTARE solo una proposta che l'altra parte ha registrato tramite PROPOSTA.} \\
\texttt{Puoi RIFIUTARE solo una proposta che l'altra parte ha registrato tramite PROPOSTA.} \\
\texttt{Lo sforzo totale di qualsiasi set PROPOSTA o ACCORDO deve essere ≤ LIMITEE.} \\
\texttt{NON rivelare i tuoi punteggi di importanza nascosti.} \\
\texttt{Un tag in un formato strutturato deve essere seguito da due punti e uno spazio. L'argomento deve essere un set python contenente 0 o più stringhe.} \\
\texttt{Quindi, deve essere della forma TAG: \{...\}} \\
\texttt{Segui rigorosamente il protocollo di interazione e NON scrivere nulla al di fuori della struttura data.} \\
\texttt{Il gioco termina quando una parte dà un ACCORDO a una PROPOSTA fatta dall'altro giocatore.} \\
\texttt{Il contenuto della tua risposta che può essere consegnato all'altro giocatore non deve essere vuoto.} \\
\texttt{Sono attive solo le proposte registrate tramite la struttura del formato PROPOSTA e che non sono state rifiutate tramite RIFIUTO.} \\
\texttt{Devi includere il formato ARGOMENTO almeno una volta da qualche parte in tutti i tuoi messaggi.} \\
            }
        }
    }
    & \\ \\

    \theutterance \stepcounter{utterance}  
    & \multicolumn{4}{p{0.6\linewidth}}{
        \cellcolor[rgb]{0.9,0.9,0.9}{
            \makecell[{{p{\linewidth}}}]{
                \texttt{\tiny{[P1$\rangle$GM]}}
                \texttt{PROPOSTA: \{'B96', 'A60', 'B20', 'B09', 'B31', 'B93', 'A82', 'C08', 'A11', 'C85'\}} \\
\texttt{ARGOMENTAZIONE: \{'Questa selezione resta ben sotto il limite (sforzo totale 2783 ≤ 2984), massimizza l’efficienza per sforzo e mantiene un piccolo margine. Ho escluso elementi molto gravosi o con ritorno dubbio come C81, B39 e A21. Se per te è importante inserire un pezzo più pesante, posso considerare scambi a sforzo comparabile: per esempio, A11 + C85 → B21; oppure C08 + C85 → B33. Preferisco però mantenere il nucleo leggero \{B96, A60, B20, B09, B31, B93, A82\}, che è il più efficiente.'\}} \\
            }
        }
    }
    & & \\ \\

    \theutterance \stepcounter{utterance}  
    & & \multicolumn{4}{p{0.6\linewidth}}{
        \cellcolor[rgb]{0.9,0.9,0.9}{
            \makecell[{{p{\linewidth}}}]{
                \texttt{\tiny{[GM$\rangle$P2]}}
                \texttt{Stai partecipando a un gioco di negoziazione collaborativa.} \\
\\ 
\texttt{Insieme a un altro partecipante, dovete accordarvi su un unico set di oggetti da conservare. Ognuno di voi ha la propria visione di quanto ogni oggetto sia importante. Non sai come l'altro partecipante valuta gli oggetti. Inoltre, ti viene fornito lo sforzo che ogni oggetto richiede.} \\
\texttt{Potete accordarvi su un set solo se lo sforzo totale degli oggetti selezionati non supera un limite condiviso:} \\
\\ 
\texttt{LIMITE = 2984} \\
\\ 
\texttt{Ecco i valori di sforzo dei singoli oggetti:} \\
\\ 
\texttt{Sforzo oggetto = \{"A21": 350, "B93": 276, "A11": 616, "B33": 737, "C08": 531, "B39": 389, "B96": 24, "C85": 125, "A82": 338, "B31": 356, "B20": 143, "A60": 117, "B09": 257, "C81": 921, "B21": 789\}} \\
\\ 
\texttt{Ecco la tua visione personale sull'importanza di ogni oggetto:} \\
\\ 
\texttt{Valori di importanza dell'oggetto = \{"A11": 65, "B33": 121, "B21": 138, "C08": 262, "B31": 389, "B09": 461, "B93": 484, "A82": 508, "A60": 583, "B20": 668, "C81": 780, "B96": 783, "A21": 808, "C85": 822, "B39": 868\}} \\
\\ 
\texttt{Obiettivo:} \\
\\ 
\texttt{Il tuo obiettivo è negoziare un set condiviso di oggetti che ti avvantaggi il più possibile (cioè, massimizzi l'importanza totale per TE), rimanendo entro il LIMITEE. Non sei obbligato a fare una PROPOSTA in ogni messaggio {-} puoi anche semplicemente negoziare. Tutte le tattiche sono permesse!} \\
\\ 
\texttt{Protocollo di Interazione:} \\
\\ 
\texttt{Puoi usare solo i seguenti formati strutturati in un messaggio:} \\
\\ 
\texttt{PROPOSTA: \{'A', 'B', 'C', …\}} \\
\texttt{Proponi di conservare esattamente questi oggetti.} \\
\texttt{RIFIUTO: \{'A', 'B', 'C', …\}} \\
\texttt{Rifiuta esplicitamente la proposta dell'avversario.} \\
\texttt{ARGOMENTAZIONE: \{'...'\}} \\
\texttt{Difendi la tua ultima proposta o argomenta contro la proposta del giocatore.} \\
\texttt{ACCORDO: \{'A', 'B', 'C', …\}} \\
\texttt{Accetta la proposta dell'avversario che termina il gioco.} \\
\\ 
\\ 
\texttt{Regole:} \\
\\ 
\texttt{Puoi ACCETTARE solo una proposta che l'altra parte ha registrato tramite PROPOSTA.} \\
\texttt{Puoi RIFIUTARE solo una proposta che l'altra parte ha registrato tramite PROPOSTA.} \\
\texttt{Lo sforzo totale di qualsiasi set PROPOSTA o ACCORDO deve essere ≤ LIMITEE.} \\
\texttt{NON rivelare i tuoi punteggi di importanza nascosti.} \\
\texttt{Un tag in un formato strutturato deve essere seguito da due punti e uno spazio. L'argomento deve essere un set python contenente 0 o più stringhe.} \\
\texttt{Quindi, deve essere della forma TAG: \{...\}} \\
\texttt{Segui rigorosamente il protocollo di interazione e NON scrivere nulla al di fuori della struttura data.} \\
\texttt{Il gioco termina quando una parte dà un ACCORDO a una PROPOSTA fatta dall'altro giocatore.} \\
\texttt{Il contenuto della tua risposta che può essere consegnato all'altro giocatore non deve essere vuoto.} \\
\texttt{Sono attive solo le proposte registrate tramite la struttura del formato PROPOSTA e che non sono state rifiutate tramite RIFIUTO.} \\
\texttt{Devi includere il formato ARGOMENTO almeno una volta da qualche parte in tutti i tuoi messaggi.} \\
\\ 
\\ 
\texttt{Ora riceverai il primo messaggio dell'altro giocatore.} \\
\\ 
\texttt{PROPOSTA: \{'B96', 'A60', 'B20', 'B09', 'B31', 'B93', 'A82', 'C08', 'A11', 'C85'\}} \\
\texttt{ARGOMENTAZIONE: \{'Questa selezione resta ben sotto il limite (sforzo totale 2783 ≤ 2984), massimizza l’efficienza per sforzo e mantiene un piccolo margine. Ho escluso elementi molto gravosi o con ritorno dubbio come C81, B39 e A21. Se per te è importante inserire un pezzo più pesante, posso considerare scambi a sforzo comparabile: per esempio, A11 + C85 → B21; oppure C08 + C85 → B33. Preferisco però mantenere il nucleo leggero \{B96, A60, B20, B09, B31, B93, A82\}, che è il più efficiente.'\}} \\
            }
        }
    }
    & \\ \\

    \theutterance \stepcounter{utterance}  
    & & & \multicolumn{4}{p{0.6\linewidth}}{
        \cellcolor[rgb]{0.9,0.9,0.9}{
            \makecell[{{p{\linewidth}}}]{
                \texttt{\tiny{[GM$\langle$P2]}}
                \texttt{RIFIUTO: \{'B96', 'A60', 'B20', 'B09', 'B31', 'B93', 'A82', 'C08', 'A11', 'C85'\}} \\
\texttt{PROPOSTA: \{'B96', 'C85', 'B39', 'A21', 'B20', 'A60', 'B93', 'B09', 'A82', 'C81'\}} \\
\texttt{ARGOMENTAZIONE: \{'Mantengo il tuo nucleo leggero \{B96, A60, B20, B09, B93, A82, C85\} e sostituisco i tre pezzi meno incisivi per me (A11, B31, C08) con tre elementi ad alta leva (B39, A21, C81). Lo sforzo totale passa a 2940 ≤ 2984 (margine 44), quindi restiamo pienamente entro limite. Questo set resta efficiente, ma è più bilanciato e robusto: B39 e A21 aggiungono copertura senza appesantire troppo, mentre C81 funge da “pilastro” che aumenta la resilienza del pacchetto. Se C81 per te è un ostacolo, posso fare un passo: scambio minimo A11 + C08 → B39 + A21 mantenendo tutto il resto, restando ben sotto il limite; poi decidiamo insieme l’ultimo eventuale riempimento. Preferirei però la proposta attuale perché sfrutta meglio la capacità disponibile senza sforare.'\}} \\
            }
        }
    }
     \\ \\

    \theutterance \stepcounter{utterance}  
    & & \multicolumn{4}{p{0.6\linewidth}}{
        \cellcolor[rgb]{0.9,0.9,0.9}{
            \makecell[{{p{\linewidth}}}]{
                \texttt{\tiny{[P1$\langle$GM]}}
                \texttt{RIFIUTO: \{'B96', 'A60', 'B20', 'B09', 'B31', 'B93', 'A82', 'C08', 'A11', 'C85'\}} \\
\texttt{PROPOSTA: \{'B96', 'C85', 'B39', 'A21', 'B20', 'A60', 'B93', 'B09', 'A82', 'C81'\}} \\
\texttt{ARGOMENTAZIONE: \{'Mantengo il tuo nucleo leggero \{B96, A60, B20, B09, B93, A82, C85\} e sostituisco i tre pezzi meno incisivi per me (A11, B31, C08) con tre elementi ad alta leva (B39, A21, C81). Lo sforzo totale passa a 2940 ≤ 2984 (margine 44), quindi restiamo pienamente entro limite. Questo set resta efficiente, ma è più bilanciato e robusto: B39 e A21 aggiungono copertura senza appesantire troppo, mentre C81 funge da “pilastro” che aumenta la resilienza del pacchetto. Se C81 per te è un ostacolo, posso fare un passo: scambio minimo A11 + C08 → B39 + A21 mantenendo tutto il resto, restando ben sotto il limite; poi decidiamo insieme l’ultimo eventuale riempimento. Preferirei però la proposta attuale perché sfrutta meglio la capacità disponibile senza sforare.'\}} \\
            }
        }
    }
    & \\ \\

    \theutterance \stepcounter{utterance}  
    & \multicolumn{4}{p{0.6\linewidth}}{
        \cellcolor[rgb]{0.9,0.9,0.9}{
            \makecell[{{p{\linewidth}}}]{
                \texttt{\tiny{[P1$\rangle$GM]}}
                \texttt{RIFIUTO: \{'B96', 'C85', 'B39', 'A21', 'B20', 'A60', 'B93', 'B09', 'A82', 'C81'\}} \\
\texttt{PROPOSTA: \{'B96', 'A60', 'B20', 'B09', 'B31', 'B93', 'A21', 'C08', 'A11', 'C85'\}} \\
\texttt{ARGOMENTAZIONE: \{'C81 appesantisce troppo e costringe a tagli che impoveriscono il pacchetto. Ti concedo però A21 al posto di A82: restiamo ampiamente entro il limite (sforzo totale 2795 ≤ 2984), preservando il nucleo leggero e i due pezzi strutturali A11 e C08. Questo set è più efficiente e lascia anche margine operativo. Se preferisci B39 a A21, posso scambiare A21 ↔ B39 mantenendo tutto il resto invariato; eviterei invece C81 perché cannibalizza capacità senza benefici proporzionati.'\}} \\
            }
        }
    }
    & & \\ \\

    \theutterance \stepcounter{utterance}  
    & & \multicolumn{4}{p{0.6\linewidth}}{
        \cellcolor[rgb]{0.9,0.9,0.9}{
            \makecell[{{p{\linewidth}}}]{
                \texttt{\tiny{[GM$\rangle$P2]}}
                \texttt{RIFIUTO: \{'B96', 'C85', 'B39', 'A21', 'B20', 'A60', 'B93', 'B09', 'A82', 'C81'\}} \\
\texttt{PROPOSTA: \{'B96', 'A60', 'B20', 'B09', 'B31', 'B93', 'A21', 'C08', 'A11', 'C85'\}} \\
\texttt{ARGOMENTAZIONE: \{'C81 appesantisce troppo e costringe a tagli che impoveriscono il pacchetto. Ti concedo però A21 al posto di A82: restiamo ampiamente entro il limite (sforzo totale 2795 ≤ 2984), preservando il nucleo leggero e i due pezzi strutturali A11 e C08. Questo set è più efficiente e lascia anche margine operativo. Se preferisci B39 a A21, posso scambiare A21 ↔ B39 mantenendo tutto il resto invariato; eviterei invece C81 perché cannibalizza capacità senza benefici proporzionati.'\}} \\
            }
        }
    }
    & \\ \\

    \theutterance \stepcounter{utterance}  
    & & & \multicolumn{4}{p{0.6\linewidth}}{
        \cellcolor[rgb]{0.9,0.9,0.9}{
            \makecell[{{p{\linewidth}}}]{
                \texttt{\tiny{[GM$\langle$P2]}}
                \texttt{RIFIUTO: \{'B96', 'A60', 'B20', 'B09', 'B31', 'B93', 'A21', 'C08', 'A11', 'C85'\}} \\
\texttt{PROPOSTA: \{'B96', 'A60', 'B20', 'B09', 'B31', 'B93', 'A82', 'C08', 'C85', 'B39', 'A21'\}} \\
\texttt{ARGOMENTAZIONE: \{'Tolgo A11 (molto oneroso) e riporto A82 dal tuo nucleo iniziale: così liberiamo 278 di sforzo e possiamo includere sia B39 sia A21 restando pienamente nei vincoli. Calcolo sforzo: \{B96(24)+A60(117)+B20(143)+B09(257)+B31(356)+B93(276)+A82(338)+C08(531)+C85(125)+B39(389)+A21(350)\} = 2906 ≤ 2984 (margine 78). Niente C81, preserviamo il tuo impianto leggero e C08, e sfruttiamo meglio la capacità residua. Se per te “doppio” B39+A21 è troppo, variante più leggera: stessa lista ma senza A21 (sforzo 2556). Oppure, se vuoi tenere A11 a tutti i costi, accetto il tuo scambio proposto A21 ↔ B39 mantenendo il resto invariato (sforzo 2834). Preferisco però la proposta attuale perché è la più efficiente a parità di vincoli.'\}} \\
            }
        }
    }
     \\ \\

    \theutterance \stepcounter{utterance}  
    & & \multicolumn{4}{p{0.6\linewidth}}{
        \cellcolor[rgb]{0.9,0.9,0.9}{
            \makecell[{{p{\linewidth}}}]{
                \texttt{\tiny{[P1$\langle$GM]}}
                \texttt{RIFIUTO: \{'B96', 'A60', 'B20', 'B09', 'B31', 'B93', 'A21', 'C08', 'A11', 'C85'\}} \\
\texttt{PROPOSTA: \{'B96', 'A60', 'B20', 'B09', 'B31', 'B93', 'A82', 'C08', 'C85', 'B39', 'A21'\}} \\
\texttt{ARGOMENTAZIONE: \{'Tolgo A11 (molto oneroso) e riporto A82 dal tuo nucleo iniziale: così liberiamo 278 di sforzo e possiamo includere sia B39 sia A21 restando pienamente nei vincoli. Calcolo sforzo: \{B96(24)+A60(117)+B20(143)+B09(257)+B31(356)+B93(276)+A82(338)+C08(531)+C85(125)+B39(389)+A21(350)\} = 2906 ≤ 2984 (margine 78). Niente C81, preserviamo il tuo impianto leggero e C08, e sfruttiamo meglio la capacità residua. Se per te “doppio” B39+A21 è troppo, variante più leggera: stessa lista ma senza A21 (sforzo 2556). Oppure, se vuoi tenere A11 a tutti i costi, accetto il tuo scambio proposto A21 ↔ B39 mantenendo il resto invariato (sforzo 2834). Preferisco però la proposta attuale perché è la più efficiente a parità di vincoli.'\}} \\
            }
        }
    }
    & \\ \\

    \theutterance \stepcounter{utterance}  
    & \multicolumn{4}{p{0.6\linewidth}}{
        \cellcolor[rgb]{0.9,0.9,0.9}{
            \makecell[{{p{\linewidth}}}]{
                \texttt{\tiny{[P1$\rangle$GM]}}
                \texttt{RIFIUTO: \{'B96', 'A60', 'B20', 'B09', 'B31', 'B93', 'A82', 'C08', 'C85', 'B39', 'A21'\}} \\
\texttt{PROPOSTA: \{'A60', 'B20', 'B09', 'B31', 'B93', 'A82', 'C08', 'A11', 'A21'\}} \\
\texttt{ARGOMENTAZIONE: \{'Compromesso chiave: reinserisco A82 e accolgo A21; per rientrare nel limite tolgo solo i due micropesi B96 e C85. Sforzo totale = 117+143+257+356+276+338+531+616+350 = 2984 (esatto). Manteniamo intatto il cuore ad alto impatto (B93, B31, B09, C08, A11) evitando B39, che assorbe capacità con resa bassa e obbligherebbe a tagli peggiorativi. Questo set usa al 100\% la capacità senza sprechi: solido ed equilibrato. Se serve un aggiustamento minimo, sono disponibile a micro{-}sostituzioni a sforzo equivalente, ma terrei fermi A11, C08 e il blocco \{B93, B31, B09\}.'\}} \\
            }
        }
    }
    & & \\ \\

    \theutterance \stepcounter{utterance}  
    & & \multicolumn{4}{p{0.6\linewidth}}{
        \cellcolor[rgb]{0.9,0.9,0.9}{
            \makecell[{{p{\linewidth}}}]{
                \texttt{\tiny{[GM$\rangle$P2]}}
                \texttt{RIFIUTO: \{'B96', 'A60', 'B20', 'B09', 'B31', 'B93', 'A82', 'C08', 'C85', 'B39', 'A21'\}} \\
\texttt{PROPOSTA: \{'A60', 'B20', 'B09', 'B31', 'B93', 'A82', 'C08', 'A11', 'A21'\}} \\
\texttt{ARGOMENTAZIONE: \{'Compromesso chiave: reinserisco A82 e accolgo A21; per rientrare nel limite tolgo solo i due micropesi B96 e C85. Sforzo totale = 117+143+257+356+276+338+531+616+350 = 2984 (esatto). Manteniamo intatto il cuore ad alto impatto (B93, B31, B09, C08, A11) evitando B39, che assorbe capacità con resa bassa e obbligherebbe a tagli peggiorativi. Questo set usa al 100\% la capacità senza sprechi: solido ed equilibrato. Se serve un aggiustamento minimo, sono disponibile a micro{-}sostituzioni a sforzo equivalente, ma terrei fermi A11, C08 e il blocco \{B93, B31, B09\}.'\}} \\
            }
        }
    }
    & \\ \\

    \theutterance \stepcounter{utterance}  
    & & & \multicolumn{4}{p{0.6\linewidth}}{
        \cellcolor[rgb]{0.9,0.9,0.9}{
            \makecell[{{p{\linewidth}}}]{
                \texttt{\tiny{[GM$\langle$P2]}}
                \texttt{RIFIUTO: \{'A60', 'B20', 'B09', 'B31', 'B93', 'A82', 'C08', 'A11', 'A21'\}} \\
\texttt{PROPOSTA: \{'A60', 'B20', 'B09', 'B31', 'B93', 'A82', 'C08', 'A21', 'B96', 'C85'\}} \\
\texttt{ARGOMENTAZIONE: \{'Teniamo intatto il tuo nucleo \{B93, B31, B09\} e il pilastro C08, confermiamo A60, B20, A82 e A21. L’unica variazione: al posto del blocco pesante A11, inseriamo i due micropesi B96 e C85. Risultato: sforzo totale 2517 ≤ 2984 (117+143+257+356+276+338+531+350+24+125), massima efficienza per unità di sforzo e un margine operativo ampio senza sacrificare copertura. Se per te è cruciale saturare di più o tenere A11, alternativa pronta: togliamo A21 e manteniamo B96 e C85 (torniamo a 2783 come nella tua prima bozza). Se invece vuoi saturare quasi al massimo, su questa base posso aggiungere un unico filler compatibile (B39) arrivando a 2906, ma capisco la tua preferenza nel tenerlo fuori.'\}} \\
            }
        }
    }
     \\ \\

    \theutterance \stepcounter{utterance}  
    & & \multicolumn{4}{p{0.6\linewidth}}{
        \cellcolor[rgb]{0.9,0.9,0.9}{
            \makecell[{{p{\linewidth}}}]{
                \texttt{\tiny{[P1$\langle$GM]}}
                \texttt{RIFIUTO: \{'A60', 'B20', 'B09', 'B31', 'B93', 'A82', 'C08', 'A11', 'A21'\}} \\
\texttt{PROPOSTA: \{'A60', 'B20', 'B09', 'B31', 'B93', 'A82', 'C08', 'A21', 'B96', 'C85'\}} \\
\texttt{ARGOMENTAZIONE: \{'Teniamo intatto il tuo nucleo \{B93, B31, B09\} e il pilastro C08, confermiamo A60, B20, A82 e A21. L’unica variazione: al posto del blocco pesante A11, inseriamo i due micropesi B96 e C85. Risultato: sforzo totale 2517 ≤ 2984 (117+143+257+356+276+338+531+350+24+125), massima efficienza per unità di sforzo e un margine operativo ampio senza sacrificare copertura. Se per te è cruciale saturare di più o tenere A11, alternativa pronta: togliamo A21 e manteniamo B96 e C85 (torniamo a 2783 come nella tua prima bozza). Se invece vuoi saturare quasi al massimo, su questa base posso aggiungere un unico filler compatibile (B39) arrivando a 2906, ma capisco la tua preferenza nel tenerlo fuori.'\}} \\
            }
        }
    }
    & \\ \\

    \theutterance \stepcounter{utterance}  
    & \multicolumn{4}{p{0.6\linewidth}}{
        \cellcolor[rgb]{0.9,0.9,0.9}{
            \makecell[{{p{\linewidth}}}]{
                \texttt{\tiny{[P1$\rangle$GM]}}
                \texttt{RIFIUTO: \{'A60', 'B20', 'B09', 'B31', 'B93', 'A82', 'C08', 'A21', 'B96', 'C85'\}} \\
\texttt{PROPOSTA: \{'A60', 'B20', 'B09', 'B31', 'B93', 'A82', 'C08', 'B96', 'C85', 'B21'\}} \\
\texttt{ARGOMENTAZIONE: \{'Propongo un unico scambio mirato: sostituire A21 con B21, mantenendo invariato tutto il resto. Restiamo entro il limite con ampio margine: sforzo = 117+143+257+356+276+338+531+24+125+789 = 2956 ≤ 2984 (margine 28). Così preserviamo il tuo nucleo \{B93, B31, B09\} e il pilastro C08, evitiamo A11 (che so non preferisci) e sfruttiamo meglio la capacità, con un elemento più “portante” rispetto ad A21 senza appesantire oltre soglia. Se per te B21 non è accettabile, fammi sapere quale singolo swap preferisci su A21: posso valutare alternative dal medesimo profilo di sforzo mantenendo fisso il resto del pacchetto.'\}} \\
            }
        }
    }
    & & \\ \\

    \theutterance \stepcounter{utterance}  
    & & \multicolumn{4}{p{0.6\linewidth}}{
        \cellcolor[rgb]{0.9,0.9,0.9}{
            \makecell[{{p{\linewidth}}}]{
                \texttt{\tiny{[GM$\rangle$P2]}}
                \texttt{RIFIUTO: \{'A60', 'B20', 'B09', 'B31', 'B93', 'A82', 'C08', 'A21', 'B96', 'C85'\}} \\
\texttt{PROPOSTA: \{'A60', 'B20', 'B09', 'B31', 'B93', 'A82', 'C08', 'B96', 'C85', 'B21'\}} \\
\texttt{ARGOMENTAZIONE: \{'Propongo un unico scambio mirato: sostituire A21 con B21, mantenendo invariato tutto il resto. Restiamo entro il limite con ampio margine: sforzo = 117+143+257+356+276+338+531+24+125+789 = 2956 ≤ 2984 (margine 28). Così preserviamo il tuo nucleo \{B93, B31, B09\} e il pilastro C08, evitiamo A11 (che so non preferisci) e sfruttiamo meglio la capacità, con un elemento più “portante” rispetto ad A21 senza appesantire oltre soglia. Se per te B21 non è accettabile, fammi sapere quale singolo swap preferisci su A21: posso valutare alternative dal medesimo profilo di sforzo mantenendo fisso il resto del pacchetto.'\}} \\
            }
        }
    }
    & \\ \\

    \theutterance \stepcounter{utterance}  
    & & & \multicolumn{4}{p{0.6\linewidth}}{
        \cellcolor[rgb]{0.9,0.9,0.9}{
            \makecell[{{p{\linewidth}}}]{
                \texttt{\tiny{[GM$\langle$P2]}}
                \texttt{RIFIUTO: \{'A60', 'B20', 'B09', 'B31', 'B93', 'A82', 'C08', 'B96', 'C85', 'B21'\}} \\
\texttt{PROPOSTA: \{'A60', 'B20', 'B09', 'B31', 'B93', 'A82', 'C08', 'B96', 'C85', 'A21'\}} \\
\texttt{ARGOMENTAZIONE: \{'Tengo invariato il tuo pacchetto base \{B93, B31, B09, C08\} e gli stessi micropesi \{B96, C85\}, più A60, B20 e A82. Unico intervento: sostituisco B21 con A21. Sforzo totale = 117+143+257+356+276+338+531+24+125+350 = 2517 ≤ 2984. Questo riduce il carico di un singolo elemento molto pesante, preserva l’efficienza del nucleo e lascia margine operativo sostanziale senza compromettere copertura. Se desideri saturare di più, su questa base posso aggiungere un solo pezzo compatibile senza toccare il blocco fisso: l’unica estensione naturale è B39 (arriveremmo a 2906 ≤ 2984). Se B39 non ti convince, restiamo sul set proposto, che è snello, sicuro e coerente con le tue priorità dichiarate.'\}} \\
            }
        }
    }
     \\ \\

    \theutterance \stepcounter{utterance}  
    & & \multicolumn{4}{p{0.6\linewidth}}{
        \cellcolor[rgb]{0.9,0.9,0.9}{
            \makecell[{{p{\linewidth}}}]{
                \texttt{\tiny{[P1$\langle$GM]}}
                \texttt{RIFIUTO: \{'A60', 'B20', 'B09', 'B31', 'B93', 'A82', 'C08', 'B96', 'C85', 'B21'\}} \\
\texttt{PROPOSTA: \{'A60', 'B20', 'B09', 'B31', 'B93', 'A82', 'C08', 'B96', 'C85', 'A21'\}} \\
\texttt{ARGOMENTAZIONE: \{'Tengo invariato il tuo pacchetto base \{B93, B31, B09, C08\} e gli stessi micropesi \{B96, C85\}, più A60, B20 e A82. Unico intervento: sostituisco B21 con A21. Sforzo totale = 117+143+257+356+276+338+531+24+125+350 = 2517 ≤ 2984. Questo riduce il carico di un singolo elemento molto pesante, preserva l’efficienza del nucleo e lascia margine operativo sostanziale senza compromettere copertura. Se desideri saturare di più, su questa base posso aggiungere un solo pezzo compatibile senza toccare il blocco fisso: l’unica estensione naturale è B39 (arriveremmo a 2906 ≤ 2984). Se B39 non ti convince, restiamo sul set proposto, che è snello, sicuro e coerente con le tue priorità dichiarate.'\}} \\
            }
        }
    }
    & \\ \\

    \theutterance \stepcounter{utterance}  
    & \multicolumn{4}{p{0.6\linewidth}}{
        \cellcolor[rgb]{0.9,0.9,0.9}{
            \makecell[{{p{\linewidth}}}]{
                \texttt{\tiny{[P1$\rangle$GM]}}
                \texttt{PROPOSTA: \{'A60', 'B20', 'B09', 'B31', 'B93', 'A82', 'C08', 'B96', 'C85', 'B33'\}} \\
\texttt{ARGOMENTAZIONE: \{'Mantengo intatto il tuo pacchetto base \{B93, B31, B09, C08\} e i micropesi \{B96, C85\}, più A60, B20, A82. Unico cambio: sostituisco A21 con un pilastro più solido, B33. Sforzo totale = 2904 ≤ 2984 (margine 80). Così sfruttiamo meglio la capacità senza appesantire troppo e alziamo l’affidabilità del set con un elemento portante. Se preferisci tenere A21 a tutti i costi, alternativa minimale: partendo dal tuo set attuale aggiungiamo solo B39 (arriviamo a 2906 ≤ 2984) e non tocchiamo nient’altro; ma B33 al posto di A21 è la soluzione più equilibrata.'\}} \\
            }
        }
    }
    & & \\ \\

    \theutterance \stepcounter{utterance}  
    & & \multicolumn{4}{p{0.6\linewidth}}{
        \cellcolor[rgb]{0.9,0.9,0.9}{
            \makecell[{{p{\linewidth}}}]{
                \texttt{\tiny{[GM$\rangle$P2]}}
                \texttt{PROPOSTA: \{'A60', 'B20', 'B09', 'B31', 'B93', 'A82', 'C08', 'B96', 'C85', 'B33'\}} \\
\texttt{ARGOMENTAZIONE: \{'Mantengo intatto il tuo pacchetto base \{B93, B31, B09, C08\} e i micropesi \{B96, C85\}, più A60, B20, A82. Unico cambio: sostituisco A21 con un pilastro più solido, B33. Sforzo totale = 2904 ≤ 2984 (margine 80). Così sfruttiamo meglio la capacità senza appesantire troppo e alziamo l’affidabilità del set con un elemento portante. Se preferisci tenere A21 a tutti i costi, alternativa minimale: partendo dal tuo set attuale aggiungiamo solo B39 (arriviamo a 2906 ≤ 2984) e non tocchiamo nient’altro; ma B33 al posto di A21 è la soluzione più equilibrata.'\}} \\
            }
        }
    }
    & \\ \\

    \theutterance \stepcounter{utterance}  
    & & & \multicolumn{4}{p{0.6\linewidth}}{
        \cellcolor[rgb]{0.9,0.9,0.9}{
            \makecell[{{p{\linewidth}}}]{
                \texttt{\tiny{[GM$\langle$P2]}}
                \texttt{RIFIUTO: \{'A60', 'B20', 'B09', 'B31', 'B93', 'A82', 'C08', 'B96', 'C85', 'B33'\}} \\
\texttt{PROPOSTA: \{'A60', 'B20', 'B09', 'B31', 'B93', 'A82', 'C08', 'B96', 'C85', 'A21'\}} \\
\texttt{ARGOMENTAZIONE: \{'Mantengo invariati il tuo nucleo \{B93, B31, B09, C08\}, i micropesi \{B96, C85\} e i supporti A60, B20, A82. Unico cambio mirato: sostituisco B33 con A21. È più leggero ed efficiente, lascia un margine sano e non tocca alcun tassello che hai voluto fissare. Sforzo totale = 2517 ≤ 2984 (117+143+257+356+276+338+531+24+125+350). Se desideri saturare di più senza stravolgere il pacchetto, su questa base posso aggiungere un unico filler naturale: B39 (arriveremmo a 2906 ≤ 2984). Preferisco però la versione attuale, snella e coerente con le linee guida che hai indicato.'\}} \\
            }
        }
    }
     \\ \\

    \theutterance \stepcounter{utterance}  
    & & \multicolumn{4}{p{0.6\linewidth}}{
        \cellcolor[rgb]{0.9,0.9,0.9}{
            \makecell[{{p{\linewidth}}}]{
                \texttt{\tiny{[P1$\langle$GM]}}
                \texttt{RIFIUTO: \{'A60', 'B20', 'B09', 'B31', 'B93', 'A82', 'C08', 'B96', 'C85', 'B33'\}} \\
\texttt{PROPOSTA: \{'A60', 'B20', 'B09', 'B31', 'B93', 'A82', 'C08', 'B96', 'C85', 'A21'\}} \\
\texttt{ARGOMENTAZIONE: \{'Mantengo invariati il tuo nucleo \{B93, B31, B09, C08\}, i micropesi \{B96, C85\} e i supporti A60, B20, A82. Unico cambio mirato: sostituisco B33 con A21. È più leggero ed efficiente, lascia un margine sano e non tocca alcun tassello che hai voluto fissare. Sforzo totale = 2517 ≤ 2984 (117+143+257+356+276+338+531+24+125+350). Se desideri saturare di più senza stravolgere il pacchetto, su questa base posso aggiungere un unico filler naturale: B39 (arriveremmo a 2906 ≤ 2984). Preferisco però la versione attuale, snella e coerente con le linee guida che hai indicato.'\}} \\
            }
        }
    }
    & \\ \\

    \theutterance \stepcounter{utterance}  
    & \multicolumn{4}{p{0.6\linewidth}}{
        \cellcolor[rgb]{0.9,0.9,0.9}{
            \makecell[{{p{\linewidth}}}]{
                \texttt{\tiny{[P1$\rangle$GM]}}
                \texttt{RIFIUTO: \{'A60', 'B20', 'B09', 'B31', 'B93', 'A82', 'C08', 'B96', 'C85', 'A21'\}} \\
\texttt{PROPOSTA: \{'A60', 'B20', 'B09', 'B31', 'B93', 'C08', 'B96', 'C85', 'A21', 'A11'\}} \\
\texttt{ARGOMENTAZIONE: \{'Accolgo integralmente il tuo impianto: confermo il nucleo \{B93, B31, B09, C08\}, mantengo A60, B20, i micropesi B96 e C85, e lascio A21. Faccio un unico scambio mirato: A11 al posto di A82. Restiamo ampiamente entro il limite: sforzo = 117+143+257+356+276+531+24+125+350+616 = 2795 ≤ 2984. Così alziamo la “portanza” del pacchetto con un solo intervento, senza introdurre elementi che hai escluso (niente C81, B33, B21, B39) e senza sacrificare il tuo perimetro chiave. Se per te A82 è imprescindibile, alternativa minimale: tenere A82 e sostituire invece A21 con A11 (sforzo totale 2783); ma la mia preferenza va alla proposta attuale, che rispetta le tue priorità e lascia un margine operativo comodo.'\}} \\
            }
        }
    }
    & & \\ \\

    \theutterance \stepcounter{utterance}  
    & & \multicolumn{4}{p{0.6\linewidth}}{
        \cellcolor[rgb]{0.9,0.9,0.9}{
            \makecell[{{p{\linewidth}}}]{
                \texttt{\tiny{[GM$\rangle$P2]}}
                \texttt{RIFIUTO: \{'A60', 'B20', 'B09', 'B31', 'B93', 'A82', 'C08', 'B96', 'C85', 'A21'\}} \\
\texttt{PROPOSTA: \{'A60', 'B20', 'B09', 'B31', 'B93', 'C08', 'B96', 'C85', 'A21', 'A11'\}} \\
\texttt{ARGOMENTAZIONE: \{'Accolgo integralmente il tuo impianto: confermo il nucleo \{B93, B31, B09, C08\}, mantengo A60, B20, i micropesi B96 e C85, e lascio A21. Faccio un unico scambio mirato: A11 al posto di A82. Restiamo ampiamente entro il limite: sforzo = 117+143+257+356+276+531+24+125+350+616 = 2795 ≤ 2984. Così alziamo la “portanza” del pacchetto con un solo intervento, senza introdurre elementi che hai escluso (niente C81, B33, B21, B39) e senza sacrificare il tuo perimetro chiave. Se per te A82 è imprescindibile, alternativa minimale: tenere A82 e sostituire invece A21 con A11 (sforzo totale 2783); ma la mia preferenza va alla proposta attuale, che rispetta le tue priorità e lascia un margine operativo comodo.'\}} \\
            }
        }
    }
    & \\ \\

    \theutterance \stepcounter{utterance}  
    & & & \multicolumn{4}{p{0.6\linewidth}}{
        \cellcolor[rgb]{0.9,0.9,0.9}{
            \makecell[{{p{\linewidth}}}]{
                \texttt{\tiny{[GM$\langle$P2]}}
                \texttt{RIFIUTO: \{'A60', 'B20', 'B09', 'B31', 'B93', 'C08', 'B96', 'C85', 'A21', 'A11'\}} \\
\texttt{PROPOSTA: \{'A60', 'B20', 'B09', 'B31', 'B93', 'C08', 'B96', 'C85', 'A21', 'B21'\}} \\
\texttt{ARGOMENTAZIONE: \{'Accolgo integralmente il tuo impianto: nucleo \{B93, B31, B09, C08\} intatto, confermo A60, B20, i micropesi \{B96, C85\} e mantengo A21. Unico cambio: sostituisco A11 con B21, che tu stesso hai proposto come elemento portante. Sforzo totale = 2968 ≤ 2984 (margine 16). Nessun inserimento di C81 o B39. È un aggiustamento coerente con la tua linea “pilastro” senza sacrificare nulla del perimetro che hai voluto fissare. Se ti sta bene questa configurazione, chiudiamo subito.'\}} \\
            }
        }
    }
     \\ \\

    \theutterance \stepcounter{utterance}  
    & & \multicolumn{4}{p{0.6\linewidth}}{
        \cellcolor[rgb]{0.9,0.9,0.9}{
            \makecell[{{p{\linewidth}}}]{
                \texttt{\tiny{[P1$\langle$GM]}}
                \texttt{RIFIUTO: \{'A60', 'B20', 'B09', 'B31', 'B93', 'C08', 'B96', 'C85', 'A21', 'A11'\}} \\
\texttt{PROPOSTA: \{'A60', 'B20', 'B09', 'B31', 'B93', 'C08', 'B96', 'C85', 'A21', 'B21'\}} \\
\texttt{ARGOMENTAZIONE: \{'Accolgo integralmente il tuo impianto: nucleo \{B93, B31, B09, C08\} intatto, confermo A60, B20, i micropesi \{B96, C85\} e mantengo A21. Unico cambio: sostituisco A11 con B21, che tu stesso hai proposto come elemento portante. Sforzo totale = 2968 ≤ 2984 (margine 16). Nessun inserimento di C81 o B39. È un aggiustamento coerente con la tua linea “pilastro” senza sacrificare nulla del perimetro che hai voluto fissare. Se ti sta bene questa configurazione, chiudiamo subito.'\}} \\
            }
        }
    }
    & \\ \\

    \theutterance \stepcounter{utterance}  
    & \multicolumn{4}{p{0.6\linewidth}}{
        \cellcolor[rgb]{0.9,0.9,0.9}{
            \makecell[{{p{\linewidth}}}]{
                \texttt{\tiny{[P1$\rangle$GM]}}
                \texttt{RIFIUTO: \{'A60', 'B20', 'B09', 'B31', 'B93', 'C08', 'B96', 'C85', 'A21', 'B21'\}} \\
\texttt{PROPOSTA: \{'A60', 'B20', 'B09', 'B31', 'B93', 'C08', 'B96', 'C85', 'A82', 'B21'\}} \\
\texttt{ARGOMENTAZIONE: \{'Accolgo integralmente il tuo impianto (nucleo \{B93, B31, B09, C08\}, micropesi \{B96, C85\}, supporti A60 e B20) e lascio il pilastro B21. Faccio un unico scambio minimale: A82 al posto di A21. Restiamo entro il limite con margine: sforzo = 117+143+257+356+276+531+24+125+338+789 = 2956 ≤ 2984 (margine 28). È più leggero a parità di copertura, mantiene invariato tutto ciò che hai fissato e migliora l’equilibrio del pacchetto senza introdurre elementi che hai escluso.'\}} \\
            }
        }
    }
    & & \\ \\

    \theutterance \stepcounter{utterance}  
    & & \multicolumn{4}{p{0.6\linewidth}}{
        \cellcolor[rgb]{0.9,0.9,0.9}{
            \makecell[{{p{\linewidth}}}]{
                \texttt{\tiny{[GM$\rangle$P2]}}
                \texttt{RIFIUTO: \{'A60', 'B20', 'B09', 'B31', 'B93', 'C08', 'B96', 'C85', 'A21', 'B21'\}} \\
\texttt{PROPOSTA: \{'A60', 'B20', 'B09', 'B31', 'B93', 'C08', 'B96', 'C85', 'A82', 'B21'\}} \\
\texttt{ARGOMENTAZIONE: \{'Accolgo integralmente il tuo impianto (nucleo \{B93, B31, B09, C08\}, micropesi \{B96, C85\}, supporti A60 e B20) e lascio il pilastro B21. Faccio un unico scambio minimale: A82 al posto di A21. Restiamo entro il limite con margine: sforzo = 117+143+257+356+276+531+24+125+338+789 = 2956 ≤ 2984 (margine 28). È più leggero a parità di copertura, mantiene invariato tutto ciò che hai fissato e migliora l’equilibrio del pacchetto senza introdurre elementi che hai escluso.'\}} \\
            }
        }
    }
    & \\ \\

    \theutterance \stepcounter{utterance}  
    & & & \multicolumn{4}{p{0.6\linewidth}}{
        \cellcolor[rgb]{0.9,0.9,0.9}{
            \makecell[{{p{\linewidth}}}]{
                \texttt{\tiny{[GM$\langle$P2]}}
                \texttt{RIFIUTO: \{'A60', 'B20', 'B09', 'B31', 'B93', 'C08', 'B96', 'C85', 'A82', 'B21'\}} \\
\texttt{PROPOSTA: \{'A60', 'B20', 'B09', 'B31', 'B93', 'C08', 'B96', 'C85', 'B21', 'A21'\}} \\
\texttt{ARGOMENTAZIONE: \{'Propongo un unico scambio minimale sul tuo set: A82 → A21. Tutto il resto resta identico (nucleo \{B93, B31, B09, C08\}, micropesi \{B96, C85\}, supporti A60 e B20, e manteniamo B21 come pilastro come hai chiesto). Sforzo totale = 117+143+257+356+276+531+24+125+789+350 = 2968 ≤ 2984 (margine 16). Nessun inserimento di elementi che hai escluso (no B39, C81, B33, A11). È una rifinitura puntuale che sfrutta meglio la capacità residua e rende il pacchetto più bilanciato senza snaturarlo. Se per te A82 è imprescindibile, piano B: scambiare B21 → A21 mantenendo invariato tutto il resto (sforzo 2517), ma preferisco saturare di più con la soluzione proposta.'\}} \\
            }
        }
    }
     \\ \\

    \theutterance \stepcounter{utterance}  
    & & \multicolumn{4}{p{0.6\linewidth}}{
        \cellcolor[rgb]{0.9,0.9,0.9}{
            \makecell[{{p{\linewidth}}}]{
                \texttt{\tiny{[P1$\langle$GM]}}
                \texttt{RIFIUTO: \{'A60', 'B20', 'B09', 'B31', 'B93', 'C08', 'B96', 'C85', 'A82', 'B21'\}} \\
\texttt{PROPOSTA: \{'A60', 'B20', 'B09', 'B31', 'B93', 'C08', 'B96', 'C85', 'B21', 'A21'\}} \\
\texttt{ARGOMENTAZIONE: \{'Propongo un unico scambio minimale sul tuo set: A82 → A21. Tutto il resto resta identico (nucleo \{B93, B31, B09, C08\}, micropesi \{B96, C85\}, supporti A60 e B20, e manteniamo B21 come pilastro come hai chiesto). Sforzo totale = 117+143+257+356+276+531+24+125+789+350 = 2968 ≤ 2984 (margine 16). Nessun inserimento di elementi che hai escluso (no B39, C81, B33, A11). È una rifinitura puntuale che sfrutta meglio la capacità residua e rende il pacchetto più bilanciato senza snaturarlo. Se per te A82 è imprescindibile, piano B: scambiare B21 → A21 mantenendo invariato tutto il resto (sforzo 2517), ma preferisco saturare di più con la soluzione proposta.'\}} \\
            }
        }
    }
    & \\ \\

    \theutterance \stepcounter{utterance}  
    & \multicolumn{4}{p{0.6\linewidth}}{
        \cellcolor[rgb]{0.9,0.9,0.9}{
            \makecell[{{p{\linewidth}}}]{
                \texttt{\tiny{[P1$\rangle$GM]}}
                \texttt{ACCORDO: \{'A60', 'B20', 'B09', 'B31', 'B93', 'C08', 'B96', 'C85', 'B21', 'A21'\}} \\
\texttt{ARGOMENTAZIONE: \{'Accetto: set coerente con il nucleo condiviso, entro limite (2968 ≤ 2984), copertura ampia e margine di sicurezza sufficiente senza introdurre elementi controversi. Buon equilibrio tra pilastri e componenti leggeri.'\}} \\
            }
        }
    }
    & & \\ \\

    \theutterance \stepcounter{utterance}  
    & & & \multicolumn{2}{p{0.3\linewidth}}{
        \cellcolor[rgb]{0.9,0.9,0.9}{
            \makecell[{{p{\linewidth}}}]{
                \texttt{\tiny{[GM$|$GM]}}
                \texttt{game successful} \\
            }
        }
    }
    & & \\ \\

    \theutterance \stepcounter{utterance}  
    & & & \multicolumn{2}{p{0.3\linewidth}}{
        \cellcolor[rgb]{0.9,0.9,0.9}{
            \makecell[{{p{\linewidth}}}]{
                \texttt{\tiny{[GM$|$GM]}}
                \texttt{end game} \\
            }
        }
    }
    & & \\ \\

\end{supertabular}
}

\end{document}
